\section{Conclusion}
\label{sec:conclusion}

We have presented a structural framework deriving time, space,
and persistent physical structure from a single axiom:
\textit{there is difference}.

The central result is the generative chain~\cref{eq:chain},
in which the causal log matrix $C_{ij}$, defined by the
self-referential equation~\cref{eq:Cij}, produces:
observational time $\tau$ (\cref{eq:tau}),
spatial distance $d_{ij}$ (\cref{eq:distance}),
log-flow $J$ (\cref{eq:flow}),
vorticity $\omega$ (\cref{eq:vorticity}),
closure density $L$ (\cref{eq:closure_density}),
and the mass spectrum of persistent structures (\cref{eq:mass}).

The system exhibits four phases (\cref{tab:phases}).
The observable universe is identified with the classical-active phase,
in which conservation laws and effective geometry emerge as stable attractors.
Physical laws are phase-dependent, not universal.

\bigskip

\begin{keybox}
\textit{The world is not given. It is constructed.
It persists because it never fully closes.}
\end{keybox}

\bigskip

\noindent
Volumes 2 and 3 develop the consequences of this framework
for the Standard Model and for gravity and cosmology, respectively.
